\scanpage{55}
\diagram{\begin{tikzcd}[row sep=small]
X^{k-i}\arrow[r]&X^k\\
{}&{[1,k']\subset[1,k]}\arrow[l,"\text{surj}"']\\
&X^I
\end{tikzcd}}
\[
I=I'\amalg I''\qquad \text{Relation d'équivalence en }I'
\]
$I$ avec une division en un ens de parties
\deleted{non vides}, comprenant la partie vide
et le choix d'une de ses parties
\[
(R,J)\qquad
\left\{\begin{array}{l}
R\subset\unsure{\mathfrak P|\mathcal P}(I)\text{ tel que }\deleted{\unsure{\Longrightarrow}}
\left\{\begin{array}{ll}
a)&\displaystyle\coprod_{\unsure{J|A}\in R}A=I\\
b)&\unsure{A,A'\in R,\ A\ne A'\ \Longrightarrow\ A\cap A'=\varnothing}\\
c)&\varnothing\in R
\end{array}\right.\\[8pt]
J\in R
\end{array}\right.
\]
\[
\varepsilon_{R,J}\varepsilon_{R',J'}
=\left\{\begin{array}{ll}0&\text{si }J\cap J'\ne\varnothing\\[6pt]&\end{array}\right.
\]
\deleted{\unsure{\S}}
Supposons $J\cap J'=\varnothing$
\[
Z_{R,J}=\left\{(x_i)_{i\in I}\text{ tels que }
\left\{\begin{array}{ll}
x_i=a&\text{si }i\in J\\
x_i=x_j&\text{si }\exists A\in R,\text{ t.q. }i,j\in A
\end{array}\right.\right\}
\]
\[
Z^a_{R,J}\cap Z^{a'}_{R',J'}=Z_{R'',J''}
\]
\margin{du pt de vue ens.}
\margin{$R''$ est la relation d'équivalence \uline{engendrée} par $R,R'$.}
\[
\begin{array}{ll}
\widetilde J&\text{la classe }\deleted{\unsure{de J}}\inserted{\text{engendrée par}}\,J\\
\widetilde J'&\text{\unsure{---}}\ \deleted{\text{engendrée par}}\,J'
\end{array}
\]
\deleted{$\widetilde J\ne\widetilde J'$}
\[
\left\{\begin{array}{l}
\widetilde J=\widetilde J'\ne\varnothing
\quad\Longrightarrow\quad
\varepsilon_{R,J}\varepsilon_{R',J'}=0\\[8pt]
\widetilde J\ne\widetilde J'\ \text{ou }\widetilde J=\widetilde J'=\varnothing
\quad\text{alors}\quad
\varepsilon_{R,J}\varepsilon_{R',J'}
=\left(\displaystyle\prod_{i\in\widetilde J\cup\widetilde J'}\unsure{y|\eta}_i\right)
\displaystyle\prod_{\substack{A''\in R''\\A''\ne\widetilde J\\A''\ne\widetilde J'}}
\varepsilon_{\unsure{R_{A''}}}^{\deleted{\unsure{A''}}}
\varepsilon_{\unsure{R'_{A''}}}^{\deleted{\unsure{A''}}}
\end{array}\right.
\]
On est ramené au cas où $J=\varnothing$, $J'=\varnothing$,
et la relation d'équivalence \uline{engendrée} par $R,R'$ est la relation
\uline{grossière}.
Alors, \unsure{ensemblistement}, on obtient la diagonale de $X^I$
\margin{54}
